% answers_short.tex —— 2009 年·普通高等学校招生考试·福建卷·理科 · 简短答案速查
% 由 tools/make_short_answers.py 自 answers.tex 的 \gksolhead/\gkanswer 提取，
% 并与 2009/2009_数学真题_福建理_答案要点.md 逐题比对后生成，勿手改（改源 MD 后重跑）。

\begin{gknotice}
\textbf{使用说明}：本节只给最终答案便于快速核对，完整推导见解析版。答案为编写者独立推导，\textbf{非官方标准答案}。
\end{gknotice}

\gkpart{选择题}
\begin{center}\begin{tabular}{ccccccccccc}
\toprule
题号 & 1 & 2 & 3 & 4 & 5 & 6 & 7 & 8 & 9 & 10 \\
\midrule
答案 & B & A & C & D & A & C & A & C & C & D \\
\bottomrule
\end{tabular}\end{center}

\gkpart{填空题}
\begin{itemize}[leftmargin=2.6em]
  \item[\textbf{11.}] $2$。
  \item[\textbf{12.}] $1$。
  \item[\textbf{13.}] $2$。
  \item[\textbf{14.}] $(-\infty,0)$。
  \item[\textbf{15.}] $5$。
\end{itemize}

\gkpart{解答题}
\begin{itemize}[leftmargin=2.6em]
  \item[\textbf{16.}] （1）$\dfrac3{31}$；（2）分布列 $\dfrac5{31},\dfrac{10}{31},\dfrac{10}{31},\dfrac5{31},\dfrac1{31}$，$E\xi=\dfrac{80}{31}$。
  \item[\textbf{17.}] （1）$\dfrac{\sqrt{10}}{10}$；（2）存在，$AS=\dfrac{\sqrt2}2$。
  \item[\textbf{18.}] （1）$A=2\sqrt3$，$\omega=\dfrac\pi6$，$|MP|=5$；（2）$N(4+\sqrt3,0)$，最长 $4+\sqrt3$。
  \item[\textbf{19.}] （1）$S(1,2\sqrt3)$ 或 $S\left(1,\dfrac{2\sqrt3}3\right)$；（2）存在，$a=\sqrt2$。
  \item[\textbf{20.}] （1）$b=2a-1$，单调区间见解析；（2）① $t=2$；② $m\in(1,3]$。
  \item[\textbf{21.}] 【A】$M^{-1}=\begin{pmatrix}-1&3\\-1&2\end{pmatrix}$，$A(2,-3)$；【B】$2$ 个；【C】$(0,2)$。
\end{itemize}
